216
13 Line and Surface Integrals
Fig. 13.2
$\eta$
$V$
$\xi_2$
$\xi_1$
$S$
$\tau_2$ $\tau_1$
$I$
precisely, we ask how to find the volume of fluid that flows across the surface $S$ per
unit time in the direction indicated by the orienting field of normals to the surface.
To solve the problem, we remark that if the velocity field of the flow is constant
and equal to $V$, then the flow per unit time across a parallelogram $\Pi$ spanned by
vectors $\xi_1$ and $\xi_2$ equals the volume of the parallelepiped constructed on the vectors
$V, \xi_1, \xi_2$. If $\eta$ is normal to $\Pi$ and we seek the flux across $\Pi$ in the direction of $\eta$,
it equals the scalar triple product $(V, \xi_1, \xi_2)$, provided $\eta$ and the frame $\xi_1, \xi_2$ give
$\Pi$ the same orientation (that is, if $(\eta, \xi_1, \xi_2)$ is a frame having the given orientation
in $\mathbb{R}^3$). If the frame $\xi_1, \xi_2$ gives the orientation opposite to the one given by $\eta$ in $\Pi$,
then the flow in the direction of $\eta$ is $- (V, \xi_1, \xi_2)$.
We now return to the original statement of the problem. For simplicity let us
assume that the entire surface $S$ admits a smooth parametrization $\varphi : I \to S \subset G$,
where $I$ is a two-dimensional interval in the plane $\mathbb{R}^2$. We partition $I$ into small
intervals $I_i$ (Fig. 13.2). We approximate the image $\varphi(I_i)$ of each such interval by
the parallelogram spanned by the images $\xi_1 = \varphi'(t_i)\tau_1$ and $\xi_2 = \varphi'(t_i)\tau_2$ of the
displacement vectors $\tau_1, \tau_2$ along the coordinate directions. Assuming that $V(x)$
varies by only a small amount inside the piece of surface $\varphi(I_i)$ and replacing $\varphi(I_i)$
by this parallelogram, we may assume that the flux $\Delta F_i$ across the piece $\varphi(I_i)$ of
the surface is equal, with small relative error, to the flux of a constant velocity field
$V(x_i) = V(\varphi(t_i))$ across the parallelogram spanned by the vectors $\xi_1, \xi_2$.
Assuming that the frame $\xi_1, \xi_2$ gives the same orientation on $S$ as $\eta$, we find
$$ \Delta F_i \approx (V(x_i), \xi_1, \xi_2). $$
Summing the elementary fluxes, we obtain
$$ F = \sum_i \Delta F_i \approx \sum_i \omega^2(x_i)(\xi_1, \xi_2), $$
where $\omega^2(x) = (V(x), \cdot, \cdot)$ is the flux 2-form (studied in Example 8 of Sect. 12.5).
If we pass to the limit, taking ever finer partitions $P$ of the interval $I$, it is natural to